\subsection{The Clifford power-multiplier obstruction}
\label{app:Clifford-power-obstruction}

Use the Clifford closure variables from \cref{sec:invariants}, with
\((m,q)=(4,3)\), \(J=|\nabla S|^2\), and \(K=\Delta S=64R\).
The instability theorem already rules out these multipliers;
the calculation below supplies the explicit wrong-sign points described
in \cref{prop:Clifford-power-obstruction}.

\begin{proof}[Construction of the wrong-sign points]
Put \(a=R^\alpha S^\gamma\). The multiplier identity in
\cref{prop:combined-power-multiplier}, together with
\cref{prop:first-order}, gives
\[
 \cL(G_{\alpha,\gamma})=a^3FQ_{\alpha,\gamma},
\]
where \(Q_{\alpha,\gamma}\) extends smoothly across \(F=0\) on
\(\{R>0,S>0\}\). We need only its restriction to
\(R=1\), \(F=0\), \(T=u^2\). Writing \(u=\lambda\),
\(v=1-\lambda\), set
\[
 \begin{aligned}
 s_\lambda&=\lambda(4-3\lambda),\\
 j_\lambda&=16\bigl(4\lambda(1-\lambda)
                      +(4-3\lambda)\lambda^2\bigr).
 \end{aligned}
\]
These are respectively \(S,J\) at the stated boundary values.
The closure table gives
\(\Gamma(F,\Gamma(F,S))=16RS+96F^2\).
Using this identity and \(K=\Delta S=64\) in
\eqref{eq:gradient-power-Q}--\eqref{eq:combined-power-Q} gives
\begin{equation}\label{eq:Clifford-boundary-coefficient}
 \begin{aligned}
 B_{\alpha,\gamma}(\lambda)
 &=48\gamma-8-\gamma(\gamma+\tfrac32)\frac{j_\lambda}{s_\lambda}\\
 &\quad+2\alpha(5-8\gamma-2\alpha)s_\lambda.
 \end{aligned}
\end{equation}
In particular,
\begin{align}
 \lim_{\lambda\downarrow0}B_{\alpha,\gamma}(\lambda)
 &=-8(\gamma-1)(2\gamma-1),
 \label{eq:obstruction-lambda0}\\
 B_{\alpha,\gamma}\left(\frac23\right)
 &=-\frac83\left[
 2\alpha^2+(8\gamma-5)\alpha+8\gamma^2-6\gamma+3
 \right].
 \label{eq:obstruction-lambda23}
\end{align}
If the required sign \(Q_{\alpha,\gamma}\geq0\) held, then
\eqref{eq:obstruction-lambda0} would force
\begin{equation}\label{eq:gamma-window}
 \frac12\leq\gamma\leq1.
\end{equation}
But the bracket in \eqref{eq:obstruction-lambda23} satisfies
\begin{align}
 &2\alpha^2+(8\gamma-5)\alpha+8\gamma^2-6\gamma+3\notag\\
 &\qquad
 =2\left(\alpha+\frac{8\gamma-5}{4}\right)^2
  +4\gamma-\frac18
 \geq\frac{15}{8}
 \quad\text{under \eqref{eq:gamma-window}}.
 \label{eq:obstruction-square}
\end{align}
Therefore the value in \eqref{eq:obstruction-lambda23} is at most
$-5$.

To realize and perturb these boundary values, put the first two matrices in
\eqref{eq:standard-first-two}.  For prescribed
$\lambda\in(0,1)$, set $u=\lambda$, $v=1-\lambda$, choose
$x\in\R^4$ with $|x|^2=u/2$, and let
\begin{equation}\label{eq:actual-y}
 y=(x,x).
\end{equation}
Then
\[
 \tau(y)=(0,u,0,\ldots,0),\qquad T=u^2.
\]
For $0<\eps<\sqrt v$, set
\[
 z_\eps=(\sqrt{v-\eps^2},\eps,0,\ldots,0).
\]
Now \(R=1\), \(T=u^2\), and
\begin{equation}\label{eq:actual-omega}
 F(y,z_\eps)=u\eps\longrightarrow0,
 \qquad
 Q_{\alpha,\gamma}(y,z_\eps)
   \longrightarrow B_{\alpha,\gamma}(\lambda).
\end{equation}
For fixed \(\alpha,\gamma\), the product rule on
\(\{F=0,R>0,S>0\}\) gives
\[
 |\nabla(R^\alpha S^\gamma F)|^2
 =R^{2\alpha}S^{2\gamma+1}>0.
\]
Hence the gradient remains nonzero for sufficiently small \(\eps\).
If \(\gamma\notin[1/2,1]\), first fix a sufficiently small
\(\lambda>0\) using \eqref{eq:obstruction-lambda0}; otherwise take
\(\lambda=2/3\).  The function \(Q_{\alpha,\gamma}\) remains negative
for all sufficiently small \(\eps>0\),
so the normalized gradient is defined and has the wrong divergence sign
at the nearby points.
\end{proof}
